%=============================================================================!
% 8.5  CUTOFFS                                                                !
%=============================================================================!
\section{Cutoffs}
\label{sec:pe-cutoffs}

A pair potential such as the Lennard-Jones energy never reaches zero as
$r$ grows, but beyond a few $\sigma$ it is very small: at $r_{\mathrm c} = 2.5\sigma$
it is $-0.0163\varepsilon$, less than 2\% of the depth of the well.
Ignoring every pair farther apart than such a cutoff distance
$r_{\mathrm c}$ leaves each atom with a fixed number of partners, on
average, at a given density, however large the system, and with the neighbour lists of
Chapter~10 makes the cost grow as $N$ in place of $N^2$. This section
compares three ways of cutting it off.

\subsection*{Three ways to cut}

The simplest choice, truncation, keeps $\varphi$ unchanged inside
$r_{\mathrm c}$ and sets it to zero outside. The energy then jumps by
$\varphi(r_{\mathrm c})$ whenever a pair crosses the cutoff, while the
force computed from \cref{eq:pe-pair-force}, which uses only $\varphi'$,
knows nothing of the jump: the force is not the slope of the energy
reported with it at $r_{\mathrm c}$, and the total energy changes by
$\mp\varphi(r_{\mathrm c})$ at every crossing, as \cref{sec:en-drift}
warned.

Truncating and shifting subtracts the constant $\varphi(r_{\mathrm c})$
inside the cutoff,
\begin{equation*}
   \varphi_{\mathrm s}(r) = \begin{cases}
      \varphi(r) - \varphi(r_{\mathrm c}) , & r < r_{\mathrm c} ,\\
      0 , & r \ge r_{\mathrm c} ,
   \end{cases}
\end{equation*}
so that the energy is continuous. A constant has no slope, so the forces
are unchanged, and the force still jumps at $r_{\mathrm c}$, from a
pull of $\varphi'(r_{\mathrm c}) = 0.039\,\varepsilon/\sigma$ for the
Lennard-Jones potential to nothing.

Switching multiplies $\varphi$ by a function $S$ that falls smoothly from
1 to 0 between a switching distance $r_{\mathrm s}$ and $r_{\mathrm c}$.
With $x = (r - r_{\mathrm s})/(r_{\mathrm c} - r_{\mathrm s})$ running
from 0 to 1 across that range,
\begin{equation}
   S(x) = 1 - 10x^3 + 15x^4 - 6x^5
   \label{eq:pe-switch}
\end{equation}
is 1 at $x = 0$ and 0 at $x = 1$ with its first and second derivatives
zero at both ends (\cref{ex:pe-switch}); these six conditions fix the six
coefficients of a polynomial of degree five, and none of lower degree
meets them all, so the energy and the force both pass smoothly to
zero. The price is a changed force between $r_{\mathrm s}$ and
$r_{\mathrm c}$. Since $S$ depends on $r$ through $x$, the chain rule gives
$\dd S/\dd r = (\dd S/\dd x)/(r_{\mathrm c} - r_{\mathrm s})$, and the
product rule then gives the force
\begin{equation*}
   -\frac{\dd(\varphi S)}{\dd r} = -\varphi'S - \frac{\varphi}{r_{\mathrm
   c} - r_{\mathrm s}}\,\frac{\dd S}{\dd x} .
\end{equation*}
There $\varphi < 0$ and $\dd S/\dd x = -30x^2(1 - x)^2 < 0$, so the second
term is negative: a pull that the original potential does not have
(\cref{fig:pe-cutoff}b).

\begin{figure}[htb]
\centering
\includegraphics{ch08_potentials/cutoff}
\caption{The Lennard-Jones potential cut at $r_{\mathrm c} = 2.5\sigma$ by
truncation, truncation and shifting, and switching between $2\sigma$ and
$2.5\sigma$, with the full potential dashed. (a) The energy; the
truncated energy jumps at $r_{\mathrm c}$. (b) The force, the same for
truncation and shifting, which jump at $r_{\mathrm c}$. (c) The total
energy of 32 atoms followed for $20\tau$ by \texttt{solve\_newton} of
\texttt{mdlab}, an adaptive solver, as its change from the start: truncated
(oxblood), and without a cutoff, shifted and switched, all on the line
at zero.}
\label{fig:pe-cutoff}
\end{figure}

\subsection*{What each does to the motion}

\Cref{fig:pe-cutoff}(c) follows a cluster of 32 Lennard-Jones atoms for
$20\tau$ with \texttt{solve\_newton} of \texttt{mdlab}, which steps
Newton's equations forward with a high-order method and chooses each step
to hold its error below a set tolerance, once with each potential, and
records the total energy every $0.025\tau$. Without a cutoff it changes
by at most $\num{1.3e-10}\,\varepsilon$, the error of the solver. With
truncation it wanders by up to $0.49\varepsilon$, in jumps of up to
$0.11\varepsilon$ between successive records, each the net effect of
several crossings (seven for the largest) of $0.0163\varepsilon$ each. The motion is that of the shifted potential,
since the forces are the same, but the energy reported with it is not
conserved. Shifted and switched, the energy is
conserved to $\num{6e-8}\,\varepsilon$, some 500 times less well than
without a cutoff: $S$ of \cref{eq:pe-switch} has a third derivative that
jumps at both ends, so a higher derivative of the force still jumps
whenever a pair crosses $r_{\mathrm s}$ or $r_{\mathrm c}$.

The shifted potential pays for its jumping force in another way. The
solver chooses its own steps to keep its error below a set tolerance,
and to do so it called for the forces 26\,465 times without a cutoff and
28\,925 times with switching, but 1\,221\,089 times with truncation or
shifting, 46 times as often. Its error, like that of any method that advances
the motion step by step, rests on the Taylor series of the motion over a
step (\cref{sec:mo-taylor}), which assumes that the force changes
smoothly; a jump in the force breaks that assumption, and the solver can
only keep its accuracy by taking tiny steps across every jump. The
fixed-step methods of Chapter~9 cannot do that, and Chapter~9 measures
what a jump in the force costs them.

Machine-learned potentials face the same choice and make the smooth one.
Behler and Parrinello multiplied each neighbour's contribution by $\frac12
[\cos(\pi r/r_{\mathrm c}) + 1]$, which falls to zero with zero slope at
the cutoff distance\cite{BehlerParrinello2007}, and TrajCast (the planned Chapter~28)
multiplies by the polynomial envelope of DimeNet\cite{Gasteiger2020}, $1 - \frac12(p + 1)(p + 2)\,x^p +
p(p + 2)\,x^{p + 1} - \frac12 p(p + 1)\,x^{p + 2}$, with $x = r/r_{\mathrm
c}$ and $p = 6$, which is zero at the cutoff with its first and second
derivatives (\cref{ex:pe-envelope}).

\begin{keyidea}
A cutoff should take both the energy and the force smoothly to zero.
Truncation makes the energy jump, so the total energy is not conserved.
Shifting makes the energy continuous but leaves the force jumping, which
an adaptive solver can follow only with far shorter steps, 46 times as
many force calls here. Switching, or a smooth envelope as in
machine-learned potentials, removes both jumps, though a jump left in a
higher derivative still costs an accurate solver some accuracy.
\end{keyidea}

\notebook{08}{cut a potential three ways and watch the energy}

\begin{selfcheck}
For argon ($\varepsilon = \SI{0.01034}{\electronvolt}$), how large is each
jump of the total energy when a pair crosses a truncated cutoff at
$2.5\sigma$?
\end{selfcheck}
